\documentclass[11pt]{article}
\usepackage[margin=1in]{geometry}
\usepackage{amsmath}
\usepackage[version=4]{mhchem}
\usepackage[T1]{fontenc}
\usepackage{lmodern}
\usepackage[authoryear,round]{natbib}
\usepackage{graphicx}
\usepackage{url}
\usepackage[hidelinks]{hyperref}
\setlength{\emergencystretch}{2em}
\newcommand{\paperfigure}[1]{\includegraphics[width=\linewidth,height=0.68\textheight,keepaspectratio]{figures/#1}}

\title{Precession-phase dependence of abrupt warming occurrence in Greenland and speleothem records}
\author{}
\date{}
\begin{document}
\maketitle

\section*{Key Points}
\begin{itemize}
\item Precession phase improves the fit to published warming-onset rates after accounting for background climate and recent events.
\item Chronology perturbations generally retain the association, but finite-event sampling leaves broad effect-size and phase uncertainty.
\item Shared information with summer insolation and dependencies in age models limit identification of a unique orbital mechanism.
\end{itemize}

\begin{abstract}
Orbital forcing modulates millennial climate variability, but its relationship to individual abrupt transitions remains uncertain. We analyze 55 published warming onsets from Greenland and a speleothem sequence centered on the penultimate glacial, using a conditional event-rate model with recent event history, benthic oxygen isotopes, and atmospheric carbon dioxide. Adding precession phase yields 0.179 bits per event of fitted information and a reduced-model bootstrap $p=0.0022$. The fitted maximum occurs near the precession-index minimum. The association generally persists under chronology perturbations, although full-model simulations reveal broad uncertainty in its strength and preferred phase. A separate 70-event synthetic Greenland catalogue gives a similar phase response and bootstrap $p=0.0138$. Summer-solstice insolation shares substantial information with precession phase. These results support phase-dependent warming occurrence within the specified models, while finite event counts and chronology dependencies constrain mechanistic interpretation.
\end{abstract}

\section*{Plain Language Summary}
Ice cores and cave deposits record repeated abrupt climate changes during past ice ages. Changes in Earth's orbit may influence the conditions in which these transitions occur, but stronger climate fluctuations do not necessarily mean that individual events follow an orbital clock. We examine the timing of 55 published warming events, allowing their occurrence rate to depend on background climate and earlier events. The description of these event times improves when we include precession, the slow change in the orientation of Earth's axis relative to its orbit. This relationship usually survives the age perturbations considered here, and a separate reconstruction shows a similar pattern. However, simulations show that the strength and precise timing of the effect remain uncertain because there are relatively few events. Summer sunlight and precession also contain overlapping information. The findings therefore suggest that some orbital configurations favored warming occurrence, without identifying a unique physical trigger or predicting the date of an individual event.

\section{Introduction}
Abrupt Greenland warmings are a defining feature of glacial climate variability. The Dansgaard--Oeschger sequence combines rapid transitions into interstadials with slower changes in the intervening states \citep{dansgaard1993evidence,rasmussen2014stratigraphic}. Its expression extends beyond Greenland: independently dated speleothems document closely associated hydroclimate changes during the last glacial period \citep{corrick2020synchronous}, while longer marine, terrestrial, and reconstructed ice-core records place millennial variability across multiple glacial cycles \citep{barker2011800,sun2021persistent,hodell20231}. Understanding when abrupt transitions become more likely therefore requires separating their relationship to slowly changing boundary conditions from variability generated within the climate system.

Several lines of evidence connect millennial variability to orbital forcing, but the statistical quantity and proxy matter. \citet{thirumalai2020methane} found precession modulation of the amplitude of millennial methane variability, without a corresponding persistent, in-phase precession modulation in Chinese speleothem oxygen-isotope variability. Their comparison cautions against treating every orbital imprint in a proxy as direct modulation of abrupt climate dynamics. Long-record analyses also reveal changes in the orbital modulation of millennial variability \citep{sun2021persistent,hodell20231}. Such amplitude patterns can reflect changes in transition magnitude or duration as well as changes in event frequency; they do not by themselves determine the probability of an individual warming onset.

Event occurrence and residence times have also been investigated directly. \citet{lohmann2018random} showed that information on both warm and cold residence times can distinguish climate-dependent behavior even when a warming sequence is compatible with a stationary random process. In Alpine speleothems from the penultimate glacial, \citet{fohlmeister2023role} found that interstadials preferentially occurred near Northern Hemisphere summer-insolation maxima. These results motivate an explicitly conditional question: does precession phase contribute information about warming occurrence after recent events and background climate are included in the same model? Ocean--sea-ice feedbacks and climate-state dependence provide possible pathways for such modulation, without requiring a deterministic external trigger \citep{dokken2013dansgaard,vettoretti2022atmospheric}. Stochastic transitions likewise remain compatible with a varying propensity for abrupt change \citep{ditlevsen2005recurrence,kleppin2015stochastic}.

Here we combine published Greenland interstadial starts with a separately dated speleothem event sequence centered on Marine Isotope Stage (MIS) 6 \citep{rasmussen2014stratigraphic,fohlmeister2023role,held2024dansgaard}. We quantify the additional fitted information supplied by precession phase, calibrate the association test by simulating event sequences under the background model, and distinguish chronology sensitivity from finite-event uncertainty in effect strength and phase. A synthetic Greenland catalogue on the speleothem-based chronology of \citet{barker2011800} provides a separate comparison. Alternative orbital predictors and model specifications test how specifically the association can be attributed to the chosen phase description. The target throughout is warming onset per unit of total observation time, rather than the triggering rate conditional on an already cold state.

\section{Data and methods}
\subsection{Event inventories and observation intervals}
The primary catalogue contains 34 North Greenland Ice Core Project (NGRIP) interstadial starts \citep{rasmussen2014stratigraphic} and 21 published speleothem transitions (Figure~\ref{fig:data_framework}; Text~S1; Table~S1). Sixteen transitions come from the Melchsee--Frutt (MF) oxygen-isotope stack in the Swiss Alps \citep{fohlmeister2023role}; five older transitions come from the Sofular carbon-isotope record in T\"urkiye \citep{held2024dansgaard}. We retain published event identities and assign reproducible numerical times near the speleothem labels using the strongest transition-direction gradient after 50-yr Gaussian smoothing. The MF--Sofular overlap guides selection of distinct events without tuning either age model (Figures~S1--S2).

Observation intervals are 12--123 kyr before 1950 CE (kyr BP) for NGRIP and 132.5--204.5 kyr BP for the speleothem sequence. These are analysis boundaries informed by source coverage, not published completeness guarantees \citep{ngrip2004high,held2024dansgaard}. The older 1.5 kyr of each interval supplies event history and is excluded from response exposure. The remaining 180 kyr contains all 55 nominal events. The intervening gap contributes neither exposure nor history. Because the speleothem interval extends outside strict MIS~6, we refer to it as MIS~6-centered.

The separate Barker catalogue uses 70 variable-threshold warming picks on the published SpeleoAge coordinate over 0--400 kyr BP, with 398.5 kyr of response exposure \citep{barker2011800}. Its 59 fixed-threshold picks provide an event-definition sensitivity. The synthetic record and its speleothem matching are distinct from direct Greenland measurements; this comparison is not a fully independent chronology replication.

\begin{figure}[p]
\centering\paperfigure{Fig01.pdf}
\caption{Event inventories and environmental context. (a) Times assigned to published events and the observation intervals adopted for NGRIP, the MF--Sofular sequence, and the separate Barker variable-threshold catalogue. The primary intervals are disconnected; their oldest 1.5 kyr is used only for history. Gray bars represent assumed catalogue observation support; thin dark lines below them mark response exposure. They do not imply equal proxy resolution or assured detection completeness. (b) Climatic precession index from La2004 \citep{laskar2004long}. (c) LR04 benthic $\delta^{18}$O \citep{lisiecki2005pliocene}. (d) Atmospheric CO$_2$ composite \citep{bereiter2015revision}. All ages use the project's BP1950 convention. Event sources are \citet{rasmussen2014stratigraphic,fohlmeister2023role,held2024dansgaard,barker2011800}.}
\label{fig:data_framework}
\end{figure}

\subsection{Covariates and conditional rate model}
Background climate is represented by LR04 benthic $\delta^{18}$O and the Antarctic atmospheric CO$_2$ composite \citep{lisiecki2005pliocene,bereiter2015revision}. LR04 combines ice-volume and deep-water-temperature influences. Orbital inputs follow La2004 \citep{laskar2004long}. We reconcile source epochs to BP1950 and hold forcing ages fixed during event-age perturbations; the precise reference year of Barker's SpeleoAge column remains an explicitly recorded assumption (Text~S2; Table~S2). Phase is interpolated between alternating precession-index extrema: minima define $0^\circ$, maxima $180^\circ$, and phase increases toward older ages.

We fit a history-dependent Poisson count model \citep{davis2003observation,truccolo2005point} in nominal 0.2-kyr bins:
\begin{align}
Y_i\mid\mathcal H_i,\mathbf x_i&\sim\operatorname{Poisson}(\Delta t_i\lambda_i),\label{eq:counts}\\
\log\lambda_i&=\beta_0+\beta_H H_i+\beta_L L_i+\beta_C C_i+\beta_S S_i
+\beta_s\sin\phi_i+\beta_c\cos\phi_i.\label{eq:rate}
\end{align}
Here $H_i$ counts events in strictly older bins whose centers fall within the preceding 1.5 kyr, $L_i$ and $C_i$ are background covariates, and $S_i$ distinguishes the speleothem and NGRIP intervals. The current bin never contributes to its own history. Actual bin duration $\Delta t_i$ enters as an exposure offset, including short boundary bins. Scalar covariates are centered and divided by their range across response bins; external predictors are evaluated at bin centers. Barker uses the same settings without $S_i$.

The background model (BG) omits the two phase terms. Both models are refitted, and their likelihood improvement is expressed as
\begin{equation}
I=\frac{\ell_{\mathrm{full}}-\ell_{\mathrm{BG}}}{N\ln2},\qquad
LR=2(\ell_{\mathrm{full}}-\ell_{\mathrm{BG}}),\label{eq:pi}
\end{equation}
where $N$ is the response event count. We call $I$ conditional predictive information (PI), in bits per event; here it measures fitted information, not validated out-of-sample skill. The fitted phase maximum is $\operatorname{atan2}(\beta_s,\beta_c)$ and the maximum/minimum phase rate ratio is $\exp[2(\beta_s^2+\beta_c^2)^{1/2}]$. Rayleigh statistics describe unadjusted event-phase concentration separately (Text~S3).

\subsection{Calibration, uncertainty, and sensitivity}
We calibrate each main association with 9,999 parametric bootstrap simulations under its fitted BG model \citep{davison1997bootstrap}. Simulation proceeds from older to younger bins, recomputing history and refitting both models; event totals vary. Empirical $p$ is $(1+\#\{LR^*\geq LR\})/10{,}000$. Other likelihood-ratio probabilities use nominal asymptotic references unless stated otherwise.

For chronology sensitivity we generate Monte Carlo (MC) age ensembles while preserving event identities and order (Texts~S2--S4). NGRIP uses cumulative Gaussian age displacements constrained by the layer-counted error scale, together with event-definition errors \citep{rasmussen2022gicc05,svensson2008sixty}. For modeled ages beyond the counted section, published uncertainty is unquantified \citep{corrick2020synchronous}; we adopt a 4.5\%-of-age nominal envelope interpreted approximately as $2\sigma$. This working scenario extends the discussion of \citet{moseley2020nalps19}; it is not a published uncertainty curve. MF uses a shared age-error factor, Sofular uses dated-depth perturbations transferred through individual age models, and nine smoothing/window combinations represent speleothem timing sensitivity. Barker uses uniform proposals within published combined control uncertainties, interpolated on SpeleoAge. Each main catalogue has 10,000 age realizations; unsupported realizations are flagged rather than clipped or replaced.

Effect precision is assessed separately for the primary catalogue (Text~S5). We compare age-only ranges (A), an approximate joint coefficient confidence region calibrated by 5,000 full-model simulations at nominal ages (B), and working ranges from 50 full-model simulations for each of 200 valid chronology generators (C). The last construction describes combined sensitivity rather than a calibrated combined confidence interval. Additional experiments vary binning, history, background-response shape, event definition, and segment phase response (Text~S6). Orbital comparisons add eccentricity, obliquity, or 65$^\circ$N summer-solstice insolation, with and without the phase block. Nine additional comparisons per catalogue form a separate nominal Holm family \citep{holm1979simple}. LR04--phase interactions are exploratory (Text~S7).

\section{Results}
\subsection{Conditional phase association and catalogue comparison}
Precession phase improves the primary event-rate fit by $I=0.1793$ bits per event ($LR=13.673$), with empirical $p=0.0022$ (Figure~\ref{fig:phase_comparison}; Figure~S8). Its nominal $p=0.00107$ is smaller, but null simulation retains evidence for an association under the specified background model. The maximum occurs at $330.0^\circ$, near the precession-index minimum, and the fitted maximum/minimum rate ratio is 4.94 with other predictors fixed. The unadjusted Rayleigh test gives $p=0.058$. These tests answer different questions: the conditional comparison incorporates climate, event history, and exposed time, whereas the descriptive phase distribution does not.

The Barker variable-threshold catalogue gives $I=0.0894$ bits per event, empirical $p=0.0138$, a maximum at $336.9^\circ$, and a rate ratio of 2.88 (Figure~S9). Thus the separately fitted response favors a similar phase sector, with a smaller point-estimated modulation. Fixed-threshold picks give $I=0.1042$, nominal $p=0.0141$, a maximum at $350.3^\circ$, and a rate ratio of 3.10 (Figure~S12). This point-age comparison indicates that the preferred sector does not depend solely on the variable threshold; it does not provide an additional independently calibrated test.

\begin{figure}[p]
\centering\paperfigure{Fig02.pdf}
\caption{Descriptive phase distributions and fitted conditional responses. Top row: primary NGRIP--speleothem catalogue; bottom row: Barker variable- and fixed-threshold catalogues (filled pink/outlined blue sectors and solid pink/dashed blue curves). (a,c) Sector radii show counts in $30^\circ$ phase sectors, without adjustment for observation exposure or background climate. Rayleigh $p$ values are 0.058, 0.079 and 0.047 for the primary, variable-threshold and fixed-threshold catalogues, respectively. (b,d) Phase multipliers $\exp(\beta_s\sin\phi+\beta_c\cos\phi)$ from separately fitted full models, holding other terms fixed. Vertical dotted lines mark fitted maxima; horizontal dotted unity lines indicate zero phase contribution to log rate, not an average event rate. PI denotes conditional predictive information, the normalized fitted likelihood gain in bits per event. Empirical association $p$ values are 0.0022 and 0.0138 for the primary and Barker variable-threshold analyses; fixed-threshold $p=0.0141$ is nominal. Minima/maxima of the precession index are $0^\circ/180^\circ$; phase increases toward older ages. Effect precision for the primary catalogue is shown in Figure~\ref{fig:effect_precision}.}
\label{fig:phase_comparison}
\end{figure}

\subsection{Chronology robustness and effect precision}
Of 10,000 joint primary age realizations, 9,982 remain within fixed observation support; 18 are flagged outside it. Among valid fits, median PI is 0.1622 bits per event, with a central 95\% working range of 0.0862--0.2175. Nominal $p<0.05$ in 98.69\% of valid fits (Figure~S6). This fraction measures sensitivity under the adopted age model, not an empirical significance level. All 10,000 Barker age fits are valid, giving median PI 0.0776 (0.0535--0.1061) and a corresponding fraction of 89.66\% (Figure~S7).

Full-model simulations reveal appreciably broader effect uncertainty than chronology perturbations alone (Figure~\ref{fig:effect_precision}). For the primary catalogue, age perturbations give a preferred-phase range of $315.9^\circ$--$343.7^\circ$ and rate-ratio range 2.92--5.96. Projecting the full-model sampling confidence region instead gives $286.5^\circ$--$374.1^\circ$ and 1.56--15.58. Angles above $360^\circ$ continue into the next cycle. The combined working experiment gives $295.9^\circ$--$365.4^\circ$ and 2.10--14.19. These constructions are not interchangeable or necessarily nested: B projects a joint confidence region, whereas C reports marginal ranges across refits. The common inference is a broadly localized phase preference with limited precision in the magnitude of modulation.

\begin{figure}[p]
\centering\paperfigure{Fig03.pdf}
\caption{Precision of the primary precession effect. A: central 95\% ranges across 9,982 valid chronology fits. B: projections of an approximate joint 95\% sine/cosine coefficient confidence region calibrated with 5,000 full-model simulations at nominal ages. C: central working ranges from 200 valid chronology generators with 50 simulations each. (a) Preferred phase, unwrapped around $330^\circ$. (b) Maximum/minimum phase rate ratio on a logarithmic axis. Dashed black lines show nominal-age estimates; dots show A/C medians and the B point estimate. (c) B joint-region boundary, every fifth C refit, observed coefficients (black dot), and zero effect (cross). (d) Fitted phase multiplier, simultaneous B envelope and pointwise C range. C includes refitting variability and is not a calibrated combined confidence interval or posterior; B and C need not be nested. Simulated event histories evolve dynamically and totals vary.}
\label{fig:effect_precision}
\end{figure}

\subsection{Model specification and orbital alternatives}
The primary association persists across the tested specifications. On a common 173-kyr response interval, all 30 bin-width/history/origin combinations yield nominal $p<0.008$, with PI 0.130--0.189 and phase maxima $323.8^\circ$--$331.4^\circ$ (Figure~S10). Alternative polynomial background responses give PI 0.171--0.180. Replacing the recent-event count with elapsed-time memory on matched support retains a phase contribution (Figure~S11; Table~S4). Allowing segment-specific phase coefficients gives nominal $p=0.808$; this does not establish equal responses, but separate coefficients are not required by this diagnostic. Sofular component and NGRIP chronology-knot alternatives likewise retain similar median PI (Table~S4).

Summer-solstice insolation shares more information with phase than do the other tested scalar orbital variables (Figure~\ref{fig:orbital_comparison}; Table~S5). Added to BG alone, insolation yields PI 0.123 for the primary catalogue and 0.052 for Barker. Once phase is included, its incremental PI falls to 0.022 and 0.001. Conversely, phase PI after insolation is 0.079 and 0.038, compared with 0.179 and 0.089 after BG alone. Neither residual phase comparison passes the separate nominal Holm correction. Eccentricity and obliquity supply little additional information after phase in either catalogue. These comparisons show overlapping statistical descriptions, rather than a unique physical driver selected by the main phase test.

\begin{figure}[p]
\centering\paperfigure{Fig04.pdf}
\caption{Orbital information conditional on different baselines. Top row: primary catalogue; bottom row: Barker variable-threshold catalogue. BG denotes an intercept, recent-event history, LR04 and CO$_2$, with an additional segment contrast for the primary catalogue; Pre denotes the sine/cosine phase block; Orb denotes eccentricity, obliquity or 65$^\circ$N summer-solstice insolation. Columns distinguish adding Orb to BG, adding Orb after Pre, and adding Pre after Orb. Circles show nominal-age PI, vertical ticks show chronology medians, and horizontal bars span the 2.5th--97.5th percentiles of valid age Monte Carlo (MC) fits (499/500 primary, 500/500 Barker). Dashed lines in the first and third columns show the nominal-age BG + Pre versus BG reference, not a significance threshold. Orb adds one coefficient and Pre two; insolation is the daily mean at the top of the atmosphere. All panels share PI units of bits per event. These are chronology sensitivity ranges, not sampling confidence intervals. The nine additional nominal tests per catalogue use Holm correction (Table~S5); their nulls differ from the main reduced-model bootstrap.}
\label{fig:orbital_comparison}
\end{figure}

\section{Discussion}
The primary result is a conditional association between precession phase and the occurrence of published warming onsets. It extends the comparison of interstadial occurrence and summer insolation in the Alpine record \citep{fohlmeister2023role} by jointly accounting for recent events, background climate, observation exposure, and uncertainty. The primary and Barker fits favor phases near the precession-index minimum, broadly corresponding to enhanced Northern Hemisphere summer insolation. This agreement is informative, but neither the fitted maximum nor its angular distance from an index extremum identifies a physical response lag.

A phase-dependent rate is compatible with stochastic transitions whose likelihood varies with the climate state. Internal variability can produce abrupt changes \citep{ditlevsen2005recurrence,kleppin2015stochastic}, while sea-ice and ocean feedbacks or CO$_2$-dependent oscillatory regimes can alter the conditions for their occurrence \citep{dokken2013dansgaard,pedro2022dansgaard,vettoretti2022atmospheric}. Orbital model experiments demonstrate sensitivity of abrupt climate dynamics to insolation configuration \citep{zhang2021direct}; they do not uniquely predict the warming-phase response fitted here. Our statistic provides a constraint on the distribution of warming times, rather than a model of the complete stadial--interstadial cycle. In particular, whole-time onset rates also depend on time spent in warm states, and cannot be equated with a cold-state triggering hazard \citep{lohmann2018random}.

The distinction between occurrence and amplitude is central to comparison with \citet{thirumalai2020methane}. Their methane--speleothem contrast concerns millennial variability envelopes, which combine proxy sensitivity with event magnitudes, durations, and frequencies. The present phase association neither overturns their Chinese-speleothem result nor demonstrates that methane amplitude is controlled by warming frequency. The records, observation intervals, and response variables differ. Together these studies motivate measuring several aspects of abrupt variability separately before assigning an orbital imprint to a particular process.

The alternative-driver results also constrain interpretation. Insolation is itself a function of multiple orbital elements, so shared information with precession is expected. A two-coefficient phase response and a one-coefficient scalar-insolation response impose different shapes; their raw PI values are not equivalent-complexity rankings. Furthermore, conditioning on LR04 and CO$_2$ can absorb pathways by which orbital forcing affects background climate. Small additional eccentricity or obliquity contributions therefore do not exclude indirect effects. Exploratory LR04--phase interactions give nominal $p=0.373$ for the primary catalogue and $p=0.0657$ for Barker (Figures~S13--S14). The latter is suggestive, but the two datasets do not establish background modulation of the phase effect.

Three limitations remain important. First, published inventories and adopted observation boundaries do not guarantee uniform detection completeness. The regional speleothem signals and the synthetic Barker reconstruction have different transfer functions, and the analysis conditions on the retained event identities. Second, uncertainty in chronology construction is only partially represented. The modeled Greenland extension inherits the ss09sea framework, whose development used a marine-isotope-stage age constraint derived from SPECMAP \citep{southon2004radiocarbon,wolff2010millennial}. This does not imply that individual NGRIP events were tuned to precession. LR04 is orbitally tuned, and Barker's chronology uses speleothem matching \citep{lisiecki2005pliocene,barker2011800}. Order-preserving age perturbations do not remove these dependencies. The MF shared factor, transferred Sofular control errors, assumed model-extension envelope and fixed forcing ages are working approximations, not complete age-model posteriors. Third, association robustness and effect precision differ. Null simulations support the specified phase test, while full-model simulations show that the available event count permits a wide range of modulation strengths. A strong point-estimated rate ratio should therefore not be treated as a precisely determined physical sensitivity.

\section{Conclusions}
Precession phase contributes fitted information about warming occurrence beyond recent events, LR04 and CO$_2$ in a 55-event Greenland--speleothem catalogue. The association generally survives the tested chronology and model variations, and a separately analyzed synthetic Greenland catalogue favors a similar phase sector. Finite-event uncertainty nonetheless leaves broad bounds on effect strength and preferred phase. Information shared with summer insolation, together with dependencies in chronology construction, prevents attribution to a unique orbital mechanism. The results support probabilistic phase organization within the specified event-rate framework and provide a quantitative target for subsequent proxy and model comparisons.

\section*{Open Research}
The Greenland event inventory follows \citet{rasmussen2014stratigraphic}; the GICC05 data are archived by \citet{rasmussen2022gicc05}. MF and Sofular source data are archived by \citet{fohlmeister2023mfData,held2024sofularData}. The synthetic Greenland event catalogue and its chronology controls are supplied by \citet{barker2011800}. Background and orbital inputs follow \citet{lisiecki2005pliocene,bereiter2015revision,laskar2004long}. Processed catalogues, observation intervals, saved model results, and reproducible figure/table scripts accompany this manuscript in the research workspace. Public code and data archive identifiers will be added before submission.

\section*{Acknowledgements and author statements}
\emph{Author names and affiliations, contributions, funding, acknowledgements, and conflict-of-interest statements await author confirmation.}

\bibliographystyle{agu}
\bibliography{reference}
\end{document}
